\clearpage
\fancyhead[L]{\small KOMET.}
\hypertarget{native-440-part-lower}{}
\begin{center}
{\large KOMET.}
\end{center}

\begin{center}
\noindent\rule{\linewidth}{0.4pt}

Göttingische gelehrte Anzeigen. Nr. 23. S. 217. 10. Februar 1821.

\noindent\rule{\linewidth}{0.4pt}
\end{center}

Der Komet, den Herr \textsc{Nicolet} am 21. Januar dieses Jahres in Paris entdeckt hatte, wurde am 30. Januar von Hofrat \textsc{Gauss} an unserer Sternwarte aufgefunden und beobachtet. Um $7^{\mathrm{u}}\,34^{\mathrm{m}}\,32^{\mathrm{s}}$ mittlerer Zeit hatte er folgende Koordinaten.

\begin{center}
\begin{tabular}{ll}
Rektaszension . . . . . . . . & $359^\circ\;27'\;7''$ \\
Nördliche Deklination . . . . . . & $16\quad4\quad36$ \\
\end{tabular}
\end{center}

Dieser Komet ist mit bloßem Auge unsichtbar, aber schon im Kometensucher gut zu erkennen. Dort zeigt sich auch der einige Grad lange Schweif. Von einem abgegrenzten Kern ist keine Spur zu sehen.
